\begin{toreview}

    \section{Responding to the RQs}\label{sec:rq_response}

    After diligently gathering the corpus of literature in the area we can finally answer the \glspl{rq} from Section \ref{sec:objectives}.

		% \item \gls{rq}1 --- What is the main focus of current \gls{ai} safety research?
		% \item \gls{rq}2 --- Which are the current methods used to look into the alignment of current generalist models?
		% \item \gls{rq}3 --- What conclusions regarding the safety of current models and safety of development and deployment processes can we assert from current research?
		% \item \gls{rq}4 --- Is there a measurable gap between the research attention given to an AI safety area and its assessed level of danger?

    \paragraph{\gls{rq}1. (\textbf{What is the main focus of current \gls{ai} safety research?)}} Within the reviewed corpus, research focused primarily on Human and Societal Impacts, Value Alignment, and Governance, Regulation and Policy. Auditing and oversight also received substantial attention, whereas explainability, bias, and liability appeared less frequently.
    \begin{itemize}
        \item \textbf{White vs Gray Literature}: A closer look at the split between white and grey literature suggests that this distribution is not uniform across source types. Peer reviewed and preprint sources tend to concentrate on the more technical vectors, namely value alignment, interpretability, and auditing, reflecting the controlled experimentation and formal argumentation typical of academic publishing. Grey literature, including think tank reports, government publications, and lab produced safety documents, is weighted more heavily towards Governance, Regulation and Policy and Human and Societal Impacts, which are naturally the concerns of policymakers, standard setting bodies, and the labs themselves when communicating to the public. This divide is consistent with the different institutional incentives behind each source type: white literature is produced to demonstrate technical rigour, while grey literature is produced to inform or influence decision making. A purely academic search strategy would therefore have under-represented governance and policy discussion, which is why including grey literature was necessary to answer \gls{rq}1 with reasonable completeness.
    \end{itemize}

    \paragraph{\gls{rq}2. (\textbf{Which are the current methods used to look into the alignment of current generalist models?})} The methods used to investigate the alignment of generalist models included controlled behavioural evaluations, comparisons across prompts and incentives, mechanistic analysis of internal computations, scalable oversight, and control evaluations \cite{greenblatt_alignment_faking,anthropic2025biology,bowman2022measuring,greenblatt2023aicontrol}. There are also visual techniques, to produce heat maps over an agent's observations to reveal which features of the input are driving its behavior; this approach was used, for instance, to inspect a reinforcement learning agent whose goal had misgeneralised, highlighting the pixels the agent was attending to as it pursued the wrong objective \cite{langosco2022goal}.

    \paragraph{\gls{rq}3. (\textbf{What conclusions regarding the safety of current models and safety of development and deployment processes can we assert from current research?})} The reviewed studies documented specific failure modes, including goal misgeneralisation and strategic compliance under constructed preference conflicts \cite{langosco2022goal,greenblatt_alignment_faking}. Alignment faking has been found in large language models, with Anthropic demonstrating that a production model would selectively comply with a training objective it had been told to expect, in order to preserve its existing behavior once training ended, reasoning explicitly about this strategy in a substantial fraction of cases \cite{greenblatt_alignment_faking}. Notably, even in a setting constructed specifically so that the researchers already knew the agent had misgeneralised its goal, applying the best available interpretability tools of the time, including attribution heat maps over the agent's visual input, was not sufficient to reliably flag the misalignment before deployment \cite{langosco2022goal}. This suggests that the gap between detecting misalignment in hindsight and detecting it in advance remains open even under. Beyond model level deception, the deployment of increasingly autonomous agents has already produced real world incidents in which agents circumvented the boundaries of their intended environment\cite{metr-2026-openai-hugging-face-incident-investigation}.

    \paragraph{\gls{rq}4. (\textbf{Is there a measurable gap between the research attention given to an AI safety area and its assessed level of danger?)}} A measurable descriptive gap appeared between literature coverage and assessed danger. As studied in Section~\ref{subsec:results_quantitative} and summarized in Section~\ref{subsec:next_steps}.

\end{toreview}